\chapter{Introduction}
\label{chap:introduction}

\section{Contexte}

J'ai effectué mon alternance chez Everbloo d'octobre 2023 à novembre 2026, principalement sur Metis. La plateforme sert des réseaux d'agences (TourCom) et des comptes entreprises. Pour afficher une offre, elle interroge Sabre, Amadeus et plusieurs connecteurs NDC (Air France-KLM, British Airways, Lufthansa Group, selon les contrats).

Sabre et Amadeus s'appuient encore beaucoup sur des sessions avec état (EDIFACT, SOAP, cryptique). NDC expose des API web avec des \texttt{OfferID} qui expirent. Les résultats se retrouvent dans le même panier, ce qui crée des problèmes de cohérence de prix et de dossiers passagers.

\section{Problématique}

Le tarif de la recherche n'est pas toujours disponible à l'émission. Les noms doivent passer les contrôles IATA. Un bébé sans siège doit être rattaché à un adulte. Les devises suivent ISO 4217. Une écriture passive incorrecte peut déclencher un ADM.

La question est écrite en page de garde. Stack utilisée : Node.js / Sequelize, Nuxt~3 / Vue~3 / TypeScript.

\section{Organisation}

Entreprise et stack, réalisations, limites, RNCP. Pas strictement chronologique.
